\section{Canonical activation and the reachable experiment}\label{sec:activation41}
No spanning hypothesis is needed on the given seeds. Their invariant
reachable space is
\begin{equation}\label{eq:reachable41}
 R_{\mathcal X}=\operatorname{span}\{hx:h\in H,\ x\in\mathcal X\},
 \qquad r=\dim R_{\mathcal X}.
\end{equation}
Orthogonality makes its orthogonal complement invariant. Every response
in \eqref{eq:law} is unchanged when its query vector is projected to
$R_{\mathcal X}$. Choose the \emph{canonical isotypic decomposition}
\[
 R_{\mathcal X}=R_0\oplus\bigoplus_{\lambda\in\Lambda}R_\lambda,
 \qquad R_0=R_{\mathcal X}^{H},
\]
where $R_\lambda$ contains all copies of the nontrivial real irreducible
type $E_\lambda$. Let $d_\lambda=\dim E_\lambda$, let $m_\lambda$ be
its real irreducible multiplicity, and put
\[
 s_\lambda=\min_{\|u\|=1,\ u\in E_\lambda}\dim(Hu),\qquad
 \sigma_{\mathcal X}=\max_{\lambda\in\Lambda}s_\lambda.
\]
For an empty set $\Lambda$ we set $\sigma_{\mathcal X}=0$.
An isometric intertwiner $J:E_\lambda\to R_\lambda$ identifies a single
irreducible copy, without choosing a decomposition into copies. The set
$\mathcal J_\lambda$ of such intertwiners is nonempty and compact: it is
the intersection of the linear intertwining equations with $J^*J=I$.

\begin{definition}[Canonical seed and query calibration]\label{def:activation41}
Let $P_\lambda$ be the isotypic orthogonal projection, and define
\begin{align}
 a_\lambda&=\max_{x\in\mathcal X,\ J\in\mathcal J_\lambda}\|J^*x\|,
 \label{eq:canonical-a41}\\
 B(J)&=\|J^*\|_{\ell^\infty_D\to\ell^2_{d_\lambda}}
      =\max_{\|u\|=1}\|Ju\|_1,\qquad
 c_\lambda=\max_{x,J}\frac{\|J^*x\|}{B(J)}.
 \label{eq:canonical-c41}
\end{align}
The $\ell^1$ and $\ell^\infty$ norms refer to the actual coordinate
queries $e_j$. Thus $c_\lambda$ is invariant under changes of model
coordinates when the queries are transformed with them, but a physical
change of the query family is not declared costless. One always has
$a_\lambda/\sqrt D\le c_\lambda\le a_\lambda$.
\end{definition}
The equality defining $B$ is norm duality. Since $J$ is an isometry,
$1\le B(J)\le\sqrt D$, so maxima exist and denominators are nonzero.

\begin{proposition}[A spectral formula in the commutant]\label{prop:activation41}
For $x\in\mathcal X$ set
\[
 C_{\lambda,x}=\int_H(hP_\lambda x)(hP_\lambda x)^*\,dm_H(h)
                       \quad\hbox{on }R_\lambda.
\]
Then
\begin{equation}\label{eq:activation-spectrum41}
 a_\lambda^2=d_\lambda\max_{x\in\mathcal X}
                         \lambda_{\max}(C_{\lambda,x}),
 \qquad
 \frac{\max_x\|P_\lambda x\|}{\sqrt{m_\lambda}}
           \le a_\lambda\le\max_x\|P_\lambda x\|.
\end{equation}
Every active isotypic type has $a_\lambda>0$. Both $a_\lambda$ and
$c_\lambda$ are independent of a splitting into equivalent copies.
For each fixed seed, $C_{\lambda,x}$ is the Hilbert--Schmidt orthogonal
projection of $(P_\lambda x)(P_\lambda x)^*$ onto the matrix commutant.
Consequently the spectral formula is finite linear algebra once the
isotypic projection is supplied.
\end{proposition}
\begin{proof}
Haar averaging is a self-adjoint idempotent on the space of real
matrices, with image the matrices commuting with $H$. This proves the
last assertion. For an isometric intertwiner $J$, the self-adjoint
operator $J^*C_{\lambda,x}J$ commutes with the irreducible action on
$E_\lambda$. Its eigenspaces are invariant, so it is scalar, even in
complex and quaternionic commutant types. Its trace is $\|J^*x\|^2$;
therefore it equals $\|J^*x\|^2I/d_\lambda$. This gives the upper
bound by $\lambda_{\max}(C_{\lambda,x})$.

Conversely, the top eigenspace of $C_{\lambda,x}$ is invariant and
contains an irreducible copy of type $\lambda$. Choose $J$ with image
in that eigenspace. The preceding scalar identity gives equality.
The trace is $\|P_\lambda x\|^2$, the ambient dimension is
$m_\lambda d_\lambda$, and the top eigenspace has dimension at least
$d_\lambda$. These facts give both norm bounds in
\eqref{eq:activation-spectrum41}. If $P_\lambda x=0$ for every seed,
then it is zero on their invariant span, contrary to the presence of
$R_\lambda$. The remaining assertions follow from the definitions.
\end{proof}

This formula is not an appeal to a preferred basis of a multiplicity
space. For instance, on two identical copies of a real-type $d$-dimensional
irreducible, a seed $(v,v)/\sqrt2$ has projection norm $1/\sqrt2$ on
each displayed copy and norm one on the diagonal copy. Its canonical
$a_\lambda$ is one. Taking a minimum over the displayed copies would
produce a coordinate-dependent calibration. This seed need not span
the ambient space; its reachable space is precisely the diagonal copy.

\begin{theorem}[Intrinsic numerical width with active seeds]
\label{thm:intrinsic-width41}
Assume $H$ is connected, $\mathcal X$ is a nonempty finite unit seed
set, and $0<\rho\le\min\{1/10,1/(8r)\}$. If
$\Lambda\ne\varnothing$, suppose there is a type $\lambda_*$ with
\begin{equation}\label{eq:dominant41}
 s_{\lambda_*}=\sigma_{\mathcal X},\qquad
 g_{E_{\lambda_*}}>0,\qquad
 0\le\varepsilon<\rho c_{\lambda_*}/2.
\end{equation}
Then
\begin{equation}\label{eq:intrinsic-width41}
 W_{N,\varepsilon}=\Theta(N^{\sigma_{\mathcal X}/2}).
\end{equation}
No action gap is required on the other types. The upper realization is
exact and uses common command rows and an $N$-dependent decoder. If
$\Lambda$ is empty, an exact constant-width realization exists.
All constants are for the fixed experiment, signal and error.
\end{theorem}
\begin{proof}
Fix a type with $g_{E_\lambda}>0$ and choose $x,J$ attaining
$c_\lambda$, writing $a=\|J^*x\|>0$ and $B=B(J)$. Track
$Y_t=U_{w_t}|_{E_\lambda}J^*x/a$ under independent commands with law
$p$, conditioned on the \emph{actual} register of any candidate machine.
For its decoder vector $d(S_N)\in[-1,1]^D$, the wordwise error implies
\[
 \E\langle JY_N,d(S_N)\rangle\ge\rho a-2\varepsilon B.
\]
Indeed $JY_N$ lies in the chosen invariant copy, its pairing with
$U_wx$ is $a$, and its $\ell^1$ norm is at most $B$ for each word.
On the other hand $\|J^*d(S_N)\|\le B$, so conditioning on $S_N$
bounds this correlation above by $Bq_N$. Thus
\begin{equation}\label{eq:kappa-canonical41}
 q_N\ge\kappa_\lambda:=\rho c_\lambda-2\varepsilon>0.
\end{equation}
Corollary~\ref{cor:action-budget41} and the all-direction orbit bound
now give $W_{N,\varepsilon}\ge cN^{s_\lambda/2}$. Apply this to
$\lambda_*$. The proof charges no free information about $J$ or the
isotypic type: these are analytical choices for a converse.

For the upper bound choose, solely to construct a machine, any orthogonal
irreducible decomposition of $R_{\mathcal X}$. Give a copy of dimension
$d_i$ the input-independent branch weight $w_i=d_i/r$. In that branch
initialize the minimal-orbit construction of Theorem~\ref{thm:orbit-upper}
at mean $(\rho/w_i)\pi_i x$. Its norm is at most
$\rho r/d_i\le1/(8d_i)$, the sufficient initialization radius there.
The same proof works for subunit initial vectors. Query $j$ uses the
restriction of the ambient functional $e_j$, of norm at most one.
The disjoint union of branch labels stores the selected branch. Its
weighted output mean is exactly $\rho U_wx$. A fixed subspace is
implemented by the scaled cross-polytope construction in the proof of
Theorem~\ref{thm:sum-width}; give it weight $\dim R_0/r$.
The resulting total label count is
\[
 O(1)+\sum_{\lambda}m_\lambda C_\lambda N^{s_\lambda/2}
                      =O(N^{\sigma_{\mathcal X}/2}).
\]
All copies are fixed in number. Command rows are common in time;
branch selection and the final two-label output cut are charged.
If only fixed vectors are reachable, retain the finite seed identity
or use the fixed polytope, and keep its state unchanged.
\end{proof}

\begin{corollary}[Canonical occupation and actual error]\label{cor:error41}
For each type with $g_{E_\lambda}>0$, set
$B_\lambda=B_{d_\lambda,A_\lambda,s_\lambda}$ from the orbit theorem.
Every machine of error $\varepsilon$ with
$\kappa_\lambda=\rho c_\lambda-2\varepsilon>0$ obeys
\[
 g_{E_\lambda}\#\{t<N:K_t\le k\}
 \le B_\lambda\kappa_\lambda^{-3-2/s_\lambda}k^{2/s_\lambda}.
\]
For a prescribed profile of peak $K$, the optimum actual error is at least
\begin{equation}\label{eq:error41}
 \frac12\max_{\lambda:g_{E_\lambda}>0}
 \left[\rho c_\lambda-
 \left(\frac{B_\lambda K^{2/s_\lambda}}
 {g_{E_\lambda}N}\right)^{1/(3+2/s_\lambda)}\right]_+.
\end{equation}
An empty maximum means zero. In particular only one calibrated expanding
type of maximal $s_\lambda$ is needed for the matched exponent.
\end{corollary}
\begin{proof}
Use \eqref{eq:kappa-canonical41} in the same occupation theorem, then
solve its peak form for $\varepsilon$ and maximize the resulting lower
bounds. If the actual error exceeds a type's threshold, its displayed
bound is already true. Compactness of all rows and decoders for a fixed
profile gives the same inequality for the optimum.
\end{proof}

\begin{proof}[Proof of Theorem~\ref{thm:main}]
Apply Theorem~\ref{thm:intrinsic-width41} and
Proposition~\ref{prop:action41}. The definitions use only the effective
command action, its reachable seed space, and the actual query norms.
\end{proof}
